\documentclass[12pt,reqno]{amsart}
\usepackage{math-review}
\providecommand{\Ric}{\operatorname{Ric}}
\providecommand{\Hess}{\operatorname{Hess}}
\providecommand{\tr}{\operatorname{tr}}
\providecommand{\Sym}{\operatorname{Sym}}

\title{From Bochner to the Lichnerowicz bound}
\date{}
\begin{document}
\maketitle

Positive Ricci curvature controls a positive Laplace eigenvalue through
one eigenfunction. The useful bridge is the Hessian trace: it retains
information that would be lost by simply discarding a nonnegative term.

\section{The eigenvalue problem}
We use the nonnegative Laplacian throughout. Gradient, covariant Hessian
and divergence integration are assumed; the precise spectral and
Bochner inputs are recalled before their use.

\begin{theorem}[Lichnerowicz]\label{thm:lich}
Let $(M^m,g)$ be a closed connected smooth Riemannian manifold of dimension
$m\geq2$. Suppose that for every $x\in M$ and $v\in T_xM$,
\[
 \Ric_g(v,v)\geq(m-1)g(v,v).
\]
For the nonnegative Laplacian $\Delta=-\operatorname{div}_g\nabla$, let
$\lambda_1$ denote its smallest positive eigenvalue on smooth real functions.
Then $\lambda_1\geq m$. This classical bound makes no equality-rigidity claim.
\end{theorem}

\subsection*{The callable inputs}
The following classical statements are imported, not proved here.
They are the exact versions required by this argument
\cite[Sections 29.1--29.2]{S-Datar}.

\begin{proposition}\label{prop:inputs}
Let $(M^m,g)$ be a closed connected smooth Riemannian manifold with $m\geq2$ and put
$\Delta=-\operatorname{div}_g\nabla$. The scalar Laplacian has a smallest
positive eigenvalue $\lambda_1>0$ with a nonzero smooth real eigenfunction
$f$ satisfying $\Delta f=\lambda_1 f$.
For every smooth real $h$ on $M$, the scalar Bochner identity in this sign convention is
\[
 -\tfrac12\Delta|\nabla h|^2
 =|\Hess h|^2-\langle\nabla\Delta h,\nabla h\rangle
   +\Ric(\nabla h,\nabla h).
\]
Here $\Hess h(X,Y)=\langle\nabla_X\nabla h,Y\rangle$, its norm is the
pointwise Hilbert--Schmidt norm, and $\tr_g\Hess h=-\Delta h$.
The spectral existence and local Bochner identity are imported classical inputs.
\end{proposition}

The source uses $L=\operatorname{div}\nabla$. Replacing $L$ by $-\Delta$
accounts for both minus signs in the displayed identity. No equality
case or classification theorem is included among the imports.

\subsection*{Integrating the identity}
For smooth $h,k$, the product rule and closedness give
\[
 0=\int_M\operatorname{div}(h\nabla k)\,dV_g
  =\int_M\bigl(\langle\nabla h,\nabla k\rangle-h\Delta k\bigr)\,dV_g.
\]
In particular, $\int_M\Delta q\,dV_g=0$ for smooth $q$. These formulas
use the Riemannian volume density and do not require an orientation.

Choose a nonzero first eigenfunction $f$ and write $\lambda=\lambda_1$.
Substitution in Bochner and integration yield
\begin{equation}\label{eq:bochner-int}
 \lambda\int_M|\nabla f|^2
 =\int_M\bigl(|\Hess f|^2+\Ric(\nabla f,\nabla f)\bigr).
\end{equation}
Here and below the measure is $dV_g$. The same integration by parts gives
\begin{equation}\label{eq:energy}
 \int_M|\nabla f|^2=\int_M f\Delta f
                    =\lambda\int_M f^2.
\end{equation}

\subsection*{Retaining the trace}
At each point choose a $g$-orthonormal frame and put
$H_{ab}=\Hess f(e_a,e_b)$. Cauchy--Schwarz on its diagonal entries gives
\[
 (\tr H)^2=\left(\sum_a H_{aa}\right)^2
 \leq m\sum_a H_{aa}^2\leq m\sum_{a,b}H_{ab}^2=m|H|^2.
\]
Since $\tr H=-\Delta f$, this supplies exactly the term needed in
\eqref{eq:bochner-int}:
\begin{equation}\label{eq:trace}
 |\Hess f|^2\geq\frac{(\Delta f)^2}{m}
                    =\frac{\lambda^2f^2}{m}.
\end{equation}

\subsection*{Closing the connection}
Let $I=\int_Mf^2>0$. Insert the curvature lower bound and
\eqref{eq:trace} into \eqref{eq:bochner-int}, then use \eqref{eq:energy}:
\[
 \lambda^2I\geq\frac{\lambda^2}{m}I+(m-1)\lambda I.
\]
Thus $\frac{m-1}{m}\lambda^2I\geq(m-1)\lambda I$. Each factor in
$(m-1)\lambda I$ is positive: $m\geq2$, $\lambda>0$, and $f\ne0$.
Division gives $\lambda\geq m$, proving Theorem~\ref{thm:lich}.

\section{Why these inputs cannot be dropped}
Discarding $|\Hess f|^2$ in \eqref{eq:bochner-int} would give only
$\lambda\geq m-1$. Its trace restores the stronger bound, which explains
why Cauchy--Schwarz is a genuine point of the proof.

Dimension one is a true boundary. On the circle of radius $R$, Ricci
curvature is zero, but its first positive eigenvalue is $R^{-2}$:
\[
 \Delta\cos\theta=R^{-2}\cos\theta.
\]
For $R>1$ this is less than $1$. Correspondingly, the factor $m-1$
in the attempted division would vanish.

On a manifold with boundary, the divergence integral acquires a boundary
term. This proof therefore establishes no boundary eigenvalue theorem.
Likewise, equality in the displayed inequalities is not itself a proof
of geometric rigidity.

\bibliographystyle{../manuscript/localrefs}
\bibliography{../generated/references}
\end{document}
